\documentclass{beamer}
\usepackage[utf8]{inputenc}

\usetheme{Berkeley}
\usecolortheme{default}

\setbeamertemplate{footline}[frame number]
\beamertemplatenavigationsymbolsempty


\title{The Peer Review Process}
\author[]{IB 514: Scientific Writing}
\date[]{}

\AtBeginSection[]
{
  \begin{frame}
    \frametitle{Overview}
    \tableofcontents[currentsection]
  \end{frame}
}

\begin{document}

% -------------------------------------------------
\begin{frame}
  \titlepage
\end{frame}

% ============================================================
\section{Pipleine}
\begin{frame}{The Peer Review Pipeline}

  \footnotesize
  \textbf{What happens after you hit submit?}
  \smallskip
  \begin{description}
    \item[\textbf{1. Submission}]
      Manuscript enters the journal's editorial system;
      author receives an acknowledgement and manuscript ID.
    \item[\textbf{2. Editorial Triage}]
      Editor-in-Chief or handling editor assesses scope fit and
      minimum quality threshold. \alert{Desk rejection} occurs 
      often within days without external review.
    \item[\textbf{3. Reviewer Invitation}]
      Editor invites 2--4 reviewers; many (most) decline. This stage
      can take days to several weeks.
    \item[\textbf{4. External Review}]
      Reviewers evaluate the manuscript and return written reports,
      within 2--8 weeks (often longer in practice).
    \item[\textbf{5. Editorial Decision}]
      Editor synthesises reviews into one of four (or more) outcomes:
      \textbf{accept}, \textbf{minor revision}, \textbf{major revision},
      or \textbf{reject} (with or without prejudice).
    \item[\textbf{6. Revision \& Resubmission}]
      Authors revise and submit response letter.
    \item[\textbf{7. Acceptance \& Production}]
      Manuscript enters copyediting, typesetting, and proofing
      before online publication.
  \end{description}

  \medskip
  \textbf{Typical timeline:} first decision in 3--5 months;
  full acceptance in 6--18 months across all revision rounds.

\end{frame}


% ------------------------------------------------------------

\section{Editor's Perspective}

\begin{frame}{The Editor's Perspective}

  \footnotesize
  \textbf{What editors are looking at}
  \begin{itemize}
    \item \textbf{Scope:} does the manuscript fit the journal's
      aims and readership?
    \item \textbf{Significance:} does it advance the field enough
      to justify publication?
    \item \textbf{Novelty:} is it sufficiently distinct from
      recent literature?
    \item \textbf{Presentation:} is the manuscript professionally
      prepared?
  \end{itemize}

  \medskip
  \textbf{Why desk rejection happens}
  \begin{itemize}
    \item Obvious scope mismatch
    \item Insufficient (perceived) methodological rigour on first read
    \item Incomplete submissions (missing figures, references, etc.)
    \item Manuscript prepared for a different journal
      (wrong format, wrong citation style)
  \end{itemize}

  \medskip
  \textbf{What this means for you as an author}
  \begin{itemize}
    \item Tailor your cover letter to the journal
    \item Polished submissions reduces editorial hesitation
    \item Desk rejection is typically not a judgement of your science ---
      it is usually a targeting problem
  \end{itemize}

\end{frame}


% ------------------------------------------------------------
\section{Reviewer's Perspective}

\begin{frame}{The Reviewer's Perspective}

  \footnotesize
  \textbf{What reviewers evaluate}
  \begin{itemize}
    \item Conceptual framing, novelty, significance of research question
    \item Appropriateness and rigour of methods
    \item Validity of results and statistical analyses
    \item Interpretation: do conclusions follow from the data?
    \item Presentation: clarity, figures, writing quality
    \item Completeness of citations and engagement with prior work
  \end{itemize}

  \medskip
  \textbf{What reviewers shouldn't be required to do}
  \begin{itemize}
    \item Editing your manuscript for you
    \item Reading a poorly organised paper charitably
  \end{itemize}

  \medskip
  \textbf{Reviewers are busy volunteers}
  \begin{itemize}
    \item Clear, well-structured manuscripts reduce reviewer
      burden and tend to produce more constructive feedback
    \item Reviewer fatigue is real: dense, difficult manuscripts
      are evaluated less generously
  \end{itemize}

  \medskip
  \textbf{Tip:} as you revise your writing, ask yourself --- \emph{am I making
  it easy for a tired expert to understand and trust this work?}

\end{frame}


% ------------------------------------------------------------
\section{Author's Perspective}
\begin{frame}{The Author's Perspective}

  \footnotesize
  \begin{description}
    \item[\textbf{Accept}]
      Rare on first submission. Minor corrections likely requested at proofing stage.
    \item[\textbf{Minor Revision}]
      The manuscript is broadly acceptable; targeted changes
      required. Typical turnaround: 2--4 weeks.
    \item[\textbf{Major Revision}]
      Substantial work required (new analyses, restructuring,
      additional experiments). Usually one more full review round.
      Typical turnaround: 4--12 weeks.
      \emph{This is the most common outcome.}
    \item[\textbf{Reject}]
      The manuscript will not be reconsidered at this journal.%
      \footnote{\scriptsize Entirely new manuscript expected for "rejection without prejudice." }
      May include encouragement to resubmit elsewhere.
  \end{description}
\end{frame}


% ------------------------------------------------------------
\subsection{First steps}
\begin{frame}{First steps}
  \textbf{First steps after receiving reviews}
  \begin{enumerate}
    \item Read once for overall tone --- then \textbf{step away}
    \item Return the next day and read for content, not emotion
    \item Identify which comments require new analysis, which
      require writing, and which require clarification only
    \item Distribute your synthesis and discuss a plan with co-authors \emph{before} revising
  \end{enumerate}
\end{frame}

% ------------------------------------------------------------

\subsection{Rejection Management}

\begin{frame}{Rejection Management}

  \footnotesize
  \textbf{Rejection is normal --- not exceptional}
  \begin{itemize}
    \item Rates can exceed 70--90\%; desk rejection of 50\%+ at high-impact journals
    \item Most published papers were rejected at least once
    \item Rejection does not mean the work is unpublishable ---
      it usually means the targeting was wrong
  \end{itemize}

  \medskip
  \textbf{What to do immediately after rejection}
  \begin{itemize}
    \item Read the decision letter and reviews carefully ---
      even a rejection often contains actionable feedback
    \item Discuss with co-authors: is the feedback
      \emph{scientific} (revise before resubmitting) or
      \emph{editorial} (wrong journal)?
    \item Identify the next target journal \emph{before}
      making any revisions
  \end{itemize}

  \medskip
  \textbf{Appealing:} 
      appropriate if review contains clear and foundational factual error
      or demonstrable misreading of the work.
      Rarely successful, can take weeks
      or months, and delays resubmission elsewhere

  \medskip
  \alert{\textbf{Rule of thumb:}} 
  Rejection sucks.  
  Get over it and aim to submit a rejected manuscript elsewhere within 4--6 weeks. 
  \textbf{Momentum matters}.

\end{frame}
% ------------------------------------------------------------



\section{Revising \& Resubmitting}

\subsection{Workflow}

\begin{frame}{Revision \& Resubmission Workflow}

  \footnotesize
  \textbf{Organise before you write}
  \begin{itemize}
    \item Extract all editor and reviewer comments into a
      single numbered list
    \item Identify overlap and dependencies: 
      \begin{itemize}
        \scriptsize
        \item Overlap across reviewers' comments
        \item Revisions requiring new analyses prior to writing 
      \end{itemize}
    \item Set a revision timeline (considering journal's request); 
    request extension early if needed
  \end{itemize}

  \medskip
  \textbf{Revise manuscript and draft response letter in parallel}
  \begin{itemize}
    \item Draft response as each comment is addressed, not retrospectively
    \item Track changes in manuscript for co-authors
    \item Keep change-tracked version for editor
  \end{itemize}

  \medskip
  \textbf{Verify}
  \begin{itemize}
    \item Every response references a line number in the revision
    \item Every comment has a response {\scriptsize (even just "\textit{Implemented on line x.}")}
    \item Changes haven't broken figure or table references
  \end{itemize}

\end{frame}



\subsection{Response Letter}

\begin{frame}{Responding to Reviewer Feedback}

  \footnotesize
  \textbf{The response document is as important as the revision}
  \begin{itemize}
    \item Editors often read your response \emph{before} re-reading
      the manuscript --- make it clear and complete
    \item Address \textbf{every} comment, even if you disagree
    \item Use a structured format: quote the comment, state your
      response, indicate the change made and its location%
      \footnote{\scriptsize See provided \LaTeX\ template.}
  \end{itemize}

  \medskip
  \textbf{Handle disagreement professionally}
  \begin{itemize}
    \item You may respectfully decline a reviewer request ---
      but you must explain why clearly and with evidence
    \item Phrase pushback as: \emph{``We respectfully disagree
      because\ldots and instead propose\ldots''}
    \item Never argue with tone; argue with data and logic
    \item If two reviewers contradict each other, flag it to
      the editor, use it to your advantage (if appropriate), and propose a resolution
  \end{itemize}

\end{frame}



% ------------------------------------------------------------

\section{Proofing Process}

\begin{frame}{The Proofing Process}

  \footnotesize
  \medskip
  Proofing is a \textbf{correction stage}, not a revision stage.\\
  Increasing done through journal's online proofing system.

  \medskip
  \textbf{Read very carefully, especially...}
  \begin{itemize}
    \item All author names, affiliations, and corresponding
      author details
    \item All equation, special character, and Greek symbol
      rendering
    \item Figures and tables in the correct order, at
      sufficient resolution, correctly captioned
  \end{itemize}

  \medskip
  \textbf{Practical notes}
  \begin{itemize}
    \item Proofing deadlines are typically 48--72 hours
    \item Read proof against manuscript line by line, not from memory
    \item Uncaught errors become permanent in published record; corrections require formal erratum
  \end{itemize}

  \begin{center}
  {\large \alert{\textbf{Remember to celebrate!}} }
  \end{center}

\end{frame}



\end{document}